\documentclass[10pt]{mypackage}

\usepackage[uprightgreeks,light]{kpfonts}
\renewcommand{\coloneq}{\coloneqq}
%\renewcommand{\mathbb}{\mathds}
\usepackage{homework}
%\usepackage{notes}

%\usepackage[ backend=bibtex, style = alphabetic, sorting=ynt ]{biblatex}
%\addbibresource{  }

\usepackage{parskip}

\fancyhf{}
\fancyhead[R]{Avinash Iyer}
\fancyhead[L]{Algebraic Topology: Homework 26}
\fancyfoot[C]{\thepage}

\setcounter{secnumdepth}{0}

\begin{document}
\RaggedRight
\section{Revised Problems}%
\begin{problem}[Homework 21, Problem 1]
  Let $\left( X,A \right)$ be a pair of spaces.
  \begin{enumerate}[(a)]
    \item Prove that $H_0\left( X,A \right) = 0$ if and only if $A$ intersects each path component of $X$ non-trivially.
    \item Prove that $H_1\left( X,A \right) = 0$ if and only if the map $H_1\left( A \right)\rightarrow H_1\left( X \right)$ is surjective and each path component of $X$ contains at most one path component of $A$.
  \end{enumerate}
\end{problem}
\begin{solution}\hfill
  \begin{enumerate}[(a)]
    \item Suppose $H_0\left( X,A \right) = 0$. From the long exact sequence, the map $i_{\ast}\colon H_0\left( A \right)\rightarrow H_0\left( X \right)$ is then surjective, meaning that the homology classes $\left[ \alpha \right]\in H_0\left( A \right)$ corresponding to a path component of $A$ generate $H_0\left( X \right)$ upon inclusion into $X$, so the collection of path components of $A$ surjects onto the path components of $X$. Thus, the intersection of $A$ with every path component of $X$ is nontrivial.

      In the reverse direction, if $A$ intersects every path component of $X$ nontrivially, then there are homology classes $\left[ \alpha \right]\in H_0\left( A \right)$ that correspond to path components of $H_0\left( X \right)$ via the inclusion of $A$ into $X$, so $i_{\ast}\colon H_0\left( A \right)\rightarrow H_0\left( X \right)$ is a surjection. In particular, this means that the map $j_{\ast}\colon H_0\left( X \right)\rightarrow H_0\left( X,A \right)$ has $\ker\left( j_{\ast} \right) = \img\left( i_{\ast} \right) = H_0\left( X \right)$, whence $H_0\left( X,A \right) = 0$.
    \item Suppose $H_1\left( X,A \right) = 0$. In the long exact sequence, the segment 
      \begin{align*}
        H_1\left( A \right)\xrightarrow{i_{\ast}} H_1\left( X \right)\xrightarrow{j_{\ast}} H_1\left( X,A \right) = 0
      \end{align*}
      implies that $i_{\ast}$ is surjective. Furthermore, the segment
      \begin{align*}
        0 = H_1\left( X,A \right)\xrightarrow{\partial}H_0\left( A \right)\xrightarrow{i_{\ast}}H_0\left( X \right)
      \end{align*}
      implies that $i_{\ast}$ is injective. If $\left[ \alpha_i \right]\in H_0\left( A \right)$ denotes the homology classes corresponding to path components of $A$, then the inclusion of these homology classes into $H_0\left( X \right)$ leads to an injection of path components of $A$ into path components of $X$.

      In the reverse direction, if every path component of $X$ contains at most one path component of $A$, then the homology classes in $H_0\left( A \right)$ corresponding to path components of $A$ inject into the homology classes of $X$ corresponding to the path components of $X$. Therefore, in the segment
      \begin{align*}
        H_1\left( X,A \right)\xrightarrow{\partial}H_0\left( A \right)\xrightarrow{i_{\ast}}H_0\left( X \right),
      \end{align*}
      we have that $i_{\ast}$ is injective, meaning that $\img\left( \partial \right) = \ker\left( i_{\ast} \right) = 0$, so we know that $\partial$ is the zero map. Similarly, if $H_1\left( A \right)\rightarrow H_1\left( X \right)$ is surjective, then in the segment
      \begin{align*}
        H_1\left( A \right)\xrightarrow{i_{\ast}}H_1\left( X \right)\xrightarrow{j_{\ast}}H_1\left( X,A \right),
      \end{align*}
      we have $\ker\left( j_{\ast} \right) = \img\left( i_{\ast} \right) = H_1\left( X \right)$, whence $j_{\ast}$ is the zero map. By exactness, since $\ker\left( \partial \right) = \img\left( j_{\ast} \right) = 0$, and $\partial$ is the zero map, it follows that $H_1\left( X,A \right) = 0$.
  \end{enumerate}
\end{solution}
\section{Current Problems}%
\begin{problem}
  Check that $j_{\ast}\circ \partial = d_q$.
\end{problem}
\begin{solution}
  Suppose $\left[ x \right]\in H_q\left( X^q,X^{q-1} \right)$ is a homology class, with $d_q\left[ x \right] = \partial \left[ x \right]\in H_q\left( X^{q-1},X^{q-2} \right)$ is the boundary homomorphism in the long exact sequence of the triple $\left( X^q,X^{q-1},X^{q-2} \right)$.

  Since both $\left( X^q,X^{q-2} \right)$ and $\left( X^{q-1},X^{q-2} \right)$ are good pairs, the quotient map $p\colon X\rightarrow X/X^{q-2}$ induces isomorphisms in (reduced) homology. For the long exact sequence of the triple $\left( X^{q},X^{q-1},X^{q-2} \right)$, this gives the following diagram, giving $p_{\ast}\partial = \partial p_{\ast}$.
  \begin{center}
    % https://tikzcd.yichuanshen.de/#N4Igdg9gJgpgziAXAbVABwnAlgFyxMJZABgBpiBdUkANwEMAbAVxiRAB12BjKCHBAL6l0mXPkIoAjOSq1GLNgAkA+gEcAFAA0AeqtI7gqgLSSBAShBCR2PASIAmGdXrNWiECsMmBW7V9P6fsb25pbCIBg24kQAzE5yrmycPHyC4ZFidihkkrIuCu7JvPxh1pkSyNK5zvJuHmq+hgIA9AbBQm3erUFGIRZWEaK2FY7VCQX1-j6dpt1eIaScdHA4-elD0ShxY-l1RamWsjBQAObwRKAAZgBOEAC2SNIgOBBIjuN77Gh013iMpSAbvdHtQXkgAKw1RLuNDKYBLFYCAFAh6Id5gxAANihE1h8PYyxwSIGKIhoNeWJxn2+vyw-wEFAEQA
\begin{tikzcd}
{H_q(X^q,X^{q-1})} \arrow[r, "\partial"] \arrow[d, "p_{\ast}"] & {H_{q-1}(X^{q-1},X^{q-2})} \arrow[d, "p_{\ast}"]  \\
{H_q(X^{q}/X^{q-2},X^{q-1}/X^{q-2})} \arrow[r, "\partial"]     & {H_{q-1}(X^{q-1}/X^{q-2},\ast)}                  
\end{tikzcd}
  \end{center}
  In the reduced homology long exact sequence of the pair $\left( X^{q-1},X^{q-2} \right)$, where we pass between reduced homology and standard homology with respect to $\ast$, this yields the following diagram.
  \begin{center}
    % https://tikzcd.yichuanshen.de/#N4Igdg9gJgpgziAXAbVABwnAlgFyxMJZABgBpiBdUkANwEMAbAVxiRAB12BjKCHBAL6l0mXPkIoAjOSq1GLNgAkA+sACOAWkkCAFAA0Aeuq1DOdODgCUIISOx4CRAEwzq9Zq0Qd2Adyyw8BlhgRQFVTW19IwihQ2MnAWtbEAx7cSIAZlc5DzZOHj5BYRTRBwkSUklZdwUvfN5+G2LUsUcpSur5TxAVY0i4mNIzCyTm0vSUFyq3LrzffxhA4NCovoEAegGNBKH2cysmu1byrOmc2u8CxoFZGCgAc3giUAAzACcIAFskaRAcCCQLnO3QAVqphjgBIcQO8vj9qP8kABWGa5LxocF7CxQ5Kw76IIGIxAANlRFwxwAhOOKeORCIBJLJoMx+xxFAEQA
\begin{tikzcd}
{H_{q-1}(X^{q-1},\ast)} \arrow[r, "j_{\ast}"] \arrow[d, "p_{\ast}"] & {\widetilde{H}_{q-1}(X^{q-1},X^{q-2})} \arrow[d, "p_{\ast}"] \\
 {H_{q-1}(X^{q-1},\ast)} \arrow[r, "j_{\ast}"]                       & {\widetilde{H}(X^{q-1}/X^{q-2},\ast)}                        
\end{tikzcd}
  \end{center}
  By commutativity, this gives that $d_q = p_{\ast}^{-1}j_{\ast}p_{\ast}\partial$, where $p_{\ast}^{-1}j_{\ast}$ emerges from the second diagram and $p_{\ast}\partial$ emerges from the first diagram and an application of commutativity. Since $p_{\ast}^{-1}j_{\ast}p_{\ast} = j_{\ast}$, this gives $j_{\ast}\partial = d_q$.
\end{solution}
\end{document}
